% archon:physics
% archon:covers ArchonPhysics/KineticRescaling.lean

\chapter{Conditional kinetic time rescaling}
\label{ch:ArchonPhysics_KineticRescaling}

This chapter concerns an effective kinetic evolution only.  The collision map
is an arbitrary time-independent map on a real normed state space, and the
equation is assumed through the named solution predicate.  Nothing here derives
that equation from the finite lattice, asserts a microscopic equilibration law,
or asserts that a threshold event is attained.

\begin{definition}\leanok
[Effective kinetic solution]
\label{def:kinetic:solves}
\lean{ArchonPhysics.KineticRescaling.SolvesKineticEquation}
For a collision map $C$, coupling $g$, initial state $x_0$, and path $D$, say
that $D$ solves the effective kinetic equation when $D(0)=x_0$ and
\[
  D'(t)=g^2 C(D(t))
\]
at every real time.
\end{definition}
\begin{proof}


\leanok
This is the definition of the conditional effective-model solution predicate.
It makes existence, uniqueness, and all properties of $C$ explicit external
hypotheses rather than consequences of the finite lattice model.
\end{proof}

\begin{theorem}\leanok
[Solution-predicate characterization]
\label{thm:kinetic:solution-iff}
\lean{ArchonPhysics.KineticRescaling.solvesKineticEquation_iff}
\uses{def:kinetic:solves}
The named predicate holds exactly when the initial-value equality and the
displayed derivative equality both hold.
\end{theorem}
\begin{proof}


\leanok
Unfold the definition.  Its two conjuncts are precisely the initial condition
and the pointwise differentiability equation.
\end{proof}

\begin{theorem}\leanok
[Rescaling a unique effective solution]
\label{thm:kinetic:solution-rescale}
\lean{ArchonPhysics.KineticRescaling.kineticSolution_rescale}
\uses{def:kinetic:solves, thm:kinetic:solution-iff}
% SOURCE: roadmap \S7, Target A (read from references/Thermal_physics_README.md)
% SOURCE QUOTE: "If a collision operator `C` is independent of `g` and
%
% ```math
% \partial_t D_g(t)=g^2 C(D_g(t)),
% ```
%
% then uniqueness implies `D_g(t)=D_1(g²t)`."
\textit{Source: roadmap \S7, Target A.}
If $D_1$ solves the equation at coupling $1$, $D_g$ solves it at coupling $g$,
and solutions at coupling $g$ with the given initial state are unique, then
\[
  D_g(t)=D_1(g^2t)
\]
for every real $t$.
\end{theorem}
% SOURCE QUOTE PROOF: No proof is supplied in the cited roadmap; the following
% is a project-local formal elaboration of its conditional implication.
\begin{proof}


\leanok
Put $E(t)=D_1(g^2t)$.  The chain rule and the coupling-one equation give
$E'(t)=g^2C(E(t))$, while $E(0)=D_1(0)$ is the stated initial state.  Thus
$E$ is another solution at coupling $g$.  The supplied uniqueness hypothesis
identifies $D_g$ with $E$.  This is conditional on the effective equation and
does not infer it from microscopic dynamics.
\end{proof}

\begin{definition}\leanok
[First state-event hitting time]
\label{def:kinetic:first-state-hit}
\lean{ArchonPhysics.KineticRescaling.firstStateHittingTime}
\uses{def:hitting:first-time}
For a state predicate $P$ and real-time trajectory $D$, define
\[
 \tau_P(D)=\inf\{t\in\overline{\mathbb R}_{\geq0}
       \mid 0<t\ \text{and}\ P(D(t.\operatorname{toReal}))\}.
\]
The extended-real convention gives $\top$ when no strictly positive event time
exists.
\end{definition}
\begin{proof}


\leanok
This is the extended-nonnegative infimum of strictly positive sampled state
event times.  It deliberately does not turn an infimum into an event time.
\end{proof}

\begin{theorem}\leanok
[First state-event infimum formula]
\label{thm:kinetic:first-state-hit-sinf}
\lean{ArchonPhysics.KineticRescaling.firstStateHittingTime_eq_sInf}
\uses{def:hitting:first-time, def:kinetic:first-state-hit}
The first state-event hitting time is exactly the displayed extended-real
infimum.
\end{theorem}
\begin{proof}


\leanok
Unfold the transparent definition.  No nonemptiness or attainment premise is
introduced by this equality.
\end{proof}

\begin{theorem}\leanok
[Transported state-event hitting times]
\label{thm:kinetic:first-state-hit-rescale}
\lean{ArchonPhysics.KineticRescaling.firstStateHittingTime_rescale}
\uses{def:kinetic:first-state-hit, thm:kinetic:first-state-hit-sinf}
% SOURCE: roadmap \S7, Target A (read from references/Thermal_physics_README.md)
% SOURCE QUOTE: "Consequently, any compatible first hitting time satisfies the
% exact identity
%
% ```math
% T_g(δ)=g^{-2}T_1(δ).
% ```"
\textit{Source: roadmap \S7, Target A.}
If $g\ne0$ and $D_g(t)=D_1(g^2t)$ at every real time, then
\[
 \tau_P(D_g)=\bigl(\operatorname{toNNReal}(g^2)\bigr)^{-1}\tau_P(D_1).
\]
\end{theorem}
% SOURCE QUOTE PROOF: No proof is supplied in the cited roadmap; the following
% is a project-local order-theoretic elaboration preserving non-attainment.
\begin{proof}


\leanok
Let $a=\operatorname{toNNReal}(g^2)$.  Since $g\ne0$, $a$ is positive and
finite, so multiplication by $a$ is an order automorphism of extended
nonnegative time, with inverse multiplication by $a^{-1}$.  The trajectory
identity transports a positive event time $t$ for $D_g$ to the positive event
time $at$ for $D_1$, and conversely.  This remains true at $\top$, whose real
projection is $0$ and which is fixed by multiplication by a positive finite
scalar.  Hence the two event-time sets are inverse scalar images.  Taking
their infima gives the formula, without choosing or assuming an attained hit.
\end{proof}
